\documentclass[10pt,a4paper]{amsart}
\usepackage[margin=19mm]{geometry}
\usepackage{amsmath,amssymb,mathtools}
\usepackage[scr=rsfs,cal=euler]{mathalfa}
\usepackage{xfrac}
\usepackage[unicode,hidelinks,bookmarks=false]{hyperref}
\newcommand{\ie}{i.e.\ }
\newcommand{\cf}{cf.\ }
\newcommand{\resp}{respectively\ }
\newcommand{\Subsection}{\subsection}
\providecommand{\nobreakdash}{\nobreak\hbox{-}}
\setlength{\emergencystretch}{2em}
\multlinegap0pt
\usepackage{amssymb}
\usepackage[normalem]{ulem}
\usepackage{cancel}
\usepackage{tikz}
\newtheorem{lemma}{Lemma}
\DeclareMathOperator{\THH}{THH}
\newcommand{\bb}[1]{\mathbb{#1}}
\newcommand{\sss}{\mathbb{S}}
\newcommand{\zz}{\mathbb{Z}}
\title{Whitehead filtrations for computations in Topological Hochschild Homology}
\author{Logan Hyslop}
\date{Source: arXiv:2311.06717v4 (CC BY 4.0)}
\begin{document}
\maketitle

\noindent\emph{Source and licence.} Whitehead filtrations for computations in Topological Hochschild Homology. Version arXiv:2311.06717v4 (CC BY 4.0), distributed by arXiv under the Creative Commons Attribution 4.0 International licence, http://creativecommons.org/licenses/by/4.0/, as recorded in the arXiv licence metadata for this exact version and shown on the abstract page https://arxiv.org/abs/2311.06717v4. Full licence notice: \texttt{licenses/arxiv-2311.06717-CC-BY-4.0.txt}.\par\smallskip

\noindent\textit{Editorial scope.} Counted unit: the article's lemma lem4.2 (source0.tex lines 914--916). The article's complete original proof of it is delivered with it (source0.tex lines 917--932, one \texttt{proof} environment of 16 lines). No mathematics is written, repaired or re-derived: the statement and the proof are the article's own lines, in the article's own notation, and no retained line is substituted or re-flowed.  Retained uncounted as the source-local context the proof needs: lines 888--890.  Nothing was omitted from any retained span, so the package carries no omission record.  No retained line contains a citation, so the wrapper carries no bibliography and no bibliography entry.  No retained line contains a figure environment, an image-inclusion command or drawing code, so no image is delivered and the package declares no asset dependency.  Editorial formatting only, in addition: an AMS \texttt{amsart} preamble that restates the article's own declarations and macros used by the retained lines, namely the environments \texttt{lemma} on separate counters, together with the standard packages those lines need.  The retained environments print in this document as Lemma 1 in order of appearance, whereas the article prints the same items with the numbering of its own full document.  The complete native source of the article as distributed in the arXiv e-print for this exact version is bundled unchanged as \texttt{sources/149886/source0.tex}.
\begin{lemma}\label{lem4.1}
	For $n>1$, we have that $$\zz_{(p)}\otimes_{\ell/v_1^n}\zz_{(p)}\simeq (\zz_{(p)}\otimes_{\ell}\zz_{(p)})\otimes_{\zz_{(p)}}(\zz_{(p)}\otimes_{\zz_{(p)}\otimes_{\ell}\ell/v_1^n}\zz_{(p)})$$ as $\bb{E}_\infty$-algebras.
\end{lemma}

\begin{lemma}\label{lem4.2}
	$\THH_*(\ell/v_1^n,\zz_{(p)})\simeq \THH_*(\ell,\zz_{(p)})\otimes_{\zz_{(p)}}\pi_*(\zz_{(p)}\otimes_{\zz_{(p)}\otimes_{\ell}\ell/v_1^n}\zz_{(p)})$.
\end{lemma}

\begin{proof}
By cyclic invariance, we have
$$\THH(\ell/v_1^n,\zz_{(p)})\simeq \THH(\zz_{(p)},\zz_{(p)}\otimes_{\ell/v_1^n}\zz_{(p)}).$$
Expanding this out and applying Lemma \ref{lem4.1}, we get
\begin{align*}
\THH(\zz_{(p)},\zz_{(p)}\otimes_{\ell/v_1^n}\zz_{(p)})&\simeq \zz_{(p)}\otimes_{\zz_{(p)}\otimes_{\sss}\zz_{(p)}}(\zz_{(p)}\otimes_{\ell/v_1^n}\zz_{(p)})\\
&\simeq \zz_{(p)}\otimes_{\zz_{(p)}\otimes_{\sss}\zz_{(p)}}((\zz_{(p)}\otimes_{\ell}\zz_{(p)})\otimes_{\zz_{(p)}}(\zz_{(p)}\otimes_{\zz_{(p)}\otimes_{\ell}\ell/v_1^n}\zz_{(p)})).
\end{align*}
Finally, by rearranging the order of the tensor products, one concludes
\begin{align*}
\THH(\ell/v_1^n,\zz_{(p)})&\simeq \zz_{(p)}\otimes_{\zz_{(p)}\otimes_{\sss}\zz_{(p)}}((\zz_{(p)}\otimes_{\ell}\zz_{(p)})\otimes_{\zz_{(p)}}(\zz_{(p)}\otimes_{\zz_{(p)}\otimes_{\ell}\ell/v_1^n}\zz_{(p)}))\\
&\simeq (\zz_{(p)}\otimes_{\zz_{(p)}\otimes_{\sss}\zz_{(p)}}(\zz_{(p)}\otimes_{\ell}\zz_{(p)}))\otimes_{\zz_{(p)}}(\zz_{(p)}\otimes_{\zz_{(p)}\otimes_{\ell}\ell/v_1^n}\zz_{(p)})\\
&\simeq \THH(\zz_{(p)},\zz_{(p)}\otimes_{\ell}\zz_{(p)})\otimes_{\zz_{(p)}}(\zz_{(p)}\otimes_{\zz_{(p)}\otimes_{\ell}\ell/v_1^n}\zz_{(p)})\\
&\simeq \THH(\ell,\zz_{(p)})\otimes_{\zz_{(p)}}(\zz_{(p)}\otimes_{\zz_{(p)}\otimes_{\ell}\ell/v_1^n}\zz_{(p)}).
\end{align*} 
\end{proof}

\end{document}
